\section{Motivation: recovery and decision-rule validation}
\label{sec:intro}

Validation of a virtual control arm must distinguish recovery of a requested
effect from reproduction of a statistic computed on the simulated patients.
These comparisons can coincide even for correctly implemented estimators.
We examine this equality in a stratified patient simulator, including
randomisation and censoring, and revisit the failed calibration that motivated
the audit.

Our motivating application separates two explanations for differentiation-marker
loss in colorectal cancer: fewer marker-producing cells (a \emph{compositional}
change), or lower output from the cells that remain (an \emph{intrinsic}
change). These explanations suggest different biological interpretations.
Single-cell measurements can distinguish them only when the relevant population
is observed in both tissues. Let $f_N,f_T$ denote the relevant cell fractions
in reference and diseased tissue, and $m_N,m_T$ their mean per-cell output.
A reference-based decomposition, related to demographic standardisation and
Oaxaca--Blinder decompositions \citep{kitagawa1955,oaxaca1973,blinder1973}, is
\begin{equation}
\underbrace{(f_T-f_N)m_N}_{\text{compositional}} +
\underbrace{f_N(m_T-m_N)}_{\text{intrinsic}} +
\underbrace{(f_T-f_N)(m_T-m_N)}_{\text{interaction}}
= f_Tm_T-f_Nm_N.
\label{eq:kitagawa}
\end{equation}
The interaction is reported separately. If no relevant cells are observed in
an arm, its sample mean is undefined; the identity cannot supply the missing
mean. In a synthetic population with no such cells, the corresponding
conditional mean itself has no referent. Substituting zero in either case
turns an unavailable contrast into an apparent effect.

Our implementation therefore reports an estimate, a wide-interval estimate, or
\notest{}, stored as \texttt{None} rather than $0.0$. The committed rule uses
$n$, the relevant-cell count in the diseased arm: estimate at $n\geq50$, flag a
wide interval at 20--49, and abstain below 20. The question is whether the
validation used to justify these cutpoints actually tested their coverage and
discrimination requirements.

\paragraph{Relation to existing validation methods.}
Generator--analysis alignment is well known. Congeniality concerns compatibility
of imputation and analysis models \citep{meng1994}; plasmode and semi-synthetic
benchmarks can favour methods sharing modelling choices
\citep{franklin2014,curth2021}; interventional evaluation offers another
comparison \citep{gentzel2019}; and generators can disadvantage good estimators
\citep{shaw2026}. The inverse crime includes reuse of forward and inverse
ingredients \citep{wirgin2004,kaipio2007}. Our contribution is an operational
audit of a chosen recovery statistic, examples showing its dependence on the
reference and trial design, and a decision-rule calibration that fails when
evaluated directly. Appendix~\ref{app:relatedwork} places the audit alongside
simulation-study methodology and selective prediction.

Simulation-based calibration (SBC) checks posterior ranks under prior-predictive
simulation \citep{talts2018}. Our point-estimator arms provide no posterior
samples, so SBC is not directly applicable. A correct posterior can pass SBC
while its point estimate equals our reference; conversely, SBC can expose
posterior errors that our residual cannot. Diagnostic sensitivity also depends
on the chosen test quantity \citep{modrak2025}. The proposed audit complements
these checks by asking whether the estimator and reference are the same
functional of the draw.
